% ==================== MARCO TEÓRICO Y TECNOLOGÍAS ====================
\section{Marco Teórico y Tecnologías}

La solución descrita en este documento se inscribe en la confluencia de tres áreas complementarias de la administración de sistemas. La virtualización permite ejecutar sistemas operativos completos como \vm{} aisladas sobre un mismo equipo físico y expone los recursos, como discos e interfaces de red, que el \host{} puede administrar de forma programática \parencite{virtualbox_vboxmanage, virtualbox_networking}. La administración de almacenamiento en GNU/Linux, por su parte, define el procedimiento que se busca automatizar: dejar un disco disponible y persistente en el sistema invitado \parencite{linux_storage_management}. Por último, la automatización mediante scripts convierte ese procedimiento manual en una operación programada, reproducible e idempotente, es decir, capaz de repetirse sin alterar el resultado ni comprometer los datos ya existentes \parencite{idempotency_concept}. Estos tres frentes convergen cuando una organización debe preparar almacenamiento en múltiples \vm{} sin intervención manual, que es precisamente el escenario abordado a continuación.

\subsection{Virtualización y Oracle VM VirtualBox}

Oracle VM VirtualBox es un hipervisor de tipo 2 (alojado, \textit{hosted}), es decir, una aplicación que se instala sobre el sistema operativo del \host{} y permite ejecutar múltiples sistemas invitados como \vm{} independientes sin requerir modificaciones al sistema anfitrión \parencite{virtualbox_vboxmanage}. Este modelo resulta apropiado para entornos de trabajo y educativos, donde la instalación sencilla y la gestión centralizada desde una única interfaz son requisitos frecuentes.

Entre las capacidades de VirtualBox relevantes para este trabajo se encuentran la gestión de discos virtuales en formatos como \texttt{.vdi} y \texttt{.vmdk}, que permiten adjuntar almacenamiento adicional a una \vm{} sin intervención sobre hardware físico, y la configuración de red en modo puente (\textit{bridged networking}), en la cual la interfaz de la \vm{} se asocia a una interfaz física del \host{} y obtiene una dirección IP propia dentro de la red local \parencite{virtualbox_networking}. Asimismo, VirtualBox expone una interfaz completa de línea de comandos, \texttt{VBoxManage}, mediante la cual es posible crear, iniciar, detener y consultar \vm{}, administrar discos y snapshots y consultar o modificar propiedades; dicho comando constituye el punto de entrada para toda la automatización descrita en este documento \parencite{virtualbox_vboxmanage}.

\subsection{Guest Additions}

Las \guestadditions{} son un conjunto de controladores de dispositivo y servicios que se instalan dentro del sistema operativo invitado y que mejoran la integración y el rendimiento entre el \host{} y la \vm{}: gráficos acelerados, portapapeles compartido, redimensionamiento dinámico de la ventana y sincronización de hora, entre otros \parencite{virtualbox_guestadditions}.

Su componente de mayor relevancia para este trabajo es el servicio de Guest Control, que habilita desde el \host{} el comando \texttt{VBoxManage guestcontrol}. A través de este mecanismo es posible ejecutar programas en el invitado, copiar archivos, crear directorios y gestionar sesiones, autenticándose con un usuario y contraseña del propio sistema invitado \parencite{virtualbox_guestcontrol, virtualbox_advanced_topics}. Para que todo ello funcione, las \guestadditions{} deben estar instaladas en la \vm{} y el servicio \texttt{VBoxService} debe encontrarse en ejecución; de ello depende la ejecución desatendida que la solución propone \parencite{virtualbox_guestadditions}.

\subsection{Comunicación Host--Guest}

Existen varios mecanismos para ejecutar comandos en el invitado desde el \host{}, cada uno con un distinto nivel de dependencia respecto de la configuración de red.

La Guest Control API (\texttt{VBoxManage guestcontrol}) permite ejecutar procesos, transferir archivos y crear directorios dentro de la \vm{} sin que el invitado exponga servicios de red adicionales, pues la comunicación la transporta el propio hipervisor por medio de las \guestadditions{} \parencite{virtualbox_guestcontrol, virtualbox_advanced_topics}. Como alternativa, SSH con autenticación por clave ofrece un canal estandarizado y ampliamente disponible en sistemas GNU/Linux, siempre que la \vm{} ejecute un servidor SSH y exista conectividad de red; su configuración comprende el intercambio de llaves y opciones de cliente como \texttt{StrictHostKeyChecking}, \texttt{UserKnownHostsFile} y \texttt{ConnectTimeout} \parencite{ssh_key_auth, ssh_config}. Finalmente, las carpetas compartidas facilitan la transferencia de archivos entre \host{} y \vm{}, pero no sustituyen la ejecución remota de comandos \parencite{virtualbox_advanced_topics}.

La solución adopta la Guest Control API por ser nativa de VirtualBox y por reducir la superficie de configuración: no exige instalar ni habilitar servicios de red en el invitado ni gestionar llaves SSH. El modo puente se configura igualmente para que la \vm{} obtenga una dirección IP propia en la red local, lo que permite identificarla y verificar su estado durante la ejecución de la automatización \parencite{virtualbox_networking, virtualbox_guestcontrol}.

\subsection{Gestión de Almacenamiento en GNU/Linux}

En GNU/Linux el almacenamiento se modela mediante dispositivos de bloques, cuyos nombres dependen del controlador: \texttt{/dev/sdX} para discos SCSI, SATA y USB, \texttt{/dev/vdX} para discos virtuales y \texttt{/dev/nvme*} para unidades NVMe \parencite{kernel_block_devices}. Sobre un disco se crean particiones con herramientas como \texttt{sfdisk} o \texttt{fdisk}, ambas de util-linux y disponibles por defecto en la mayoría de distribuciones; la tabla puede ser MBR o GPT, y esta última es más moderna y recomendada \parencite{linux_storage_management}.

Cada partición requiere un sistema de archivos, por ejemplo ext4 o xfs, que se crea con la utilidad correspondiente (\texttt{mkfs.ext4}, \texttt{mkfs.xfs}) \parencite{linux_storage_management}. Dicho sistema de archivos se hace accesible en un punto de montaje de la jerarquía de directorios, para la cual el estándar de jerarquía del sistema de archivos reserva rutas como \texttt{/mnt} y \texttt{/media} \parencite{lsb_fhs}. La operación \texttt{mount} asocia la partición con su punto de montaje, pero ese vínculo no persiste tras un reinicio salvo que se declare en \texttt{/etc/fstab} o en una unidad de montaje de systemd \parencite{linux_storage_management, systemd_mount}.

Para que la identificación sea estable, las entradas de \texttt{/etc/fstab} deben referenciar el sistema de archivos mediante su UUID, obtenido con \texttt{blkid}, y no mediante \texttt{/dev/sdX}, ya que la letra asignada a un dispositivo puede cambiar entre reinicios o al conectar otro disco \parencite{uuid_fstab}. Una vez aplicada la configuración, su verificación se realiza con herramientas como \texttt{lsblk}, \texttt{findmnt} y \texttt{df -h} \parencite{linux_storage_management}.

\subsection{Automatización con Bash}
\label{sub:bash}

La automatización de la preparación de almacenamiento exige un lenguaje disponible de forma nativa en el \host{}, con acceso directo a las utilidades de línea de comandos de GNU/Linux y con la madurez necesaria para construir scripts robustos. La solución se implementa íntegramente en Bash, intérprete de comandos presente de forma predeterminada en los sistemas GNU/Linux y macOS y heredero del lenguaje de comandos definido por la norma POSIX, lo que favorece su portabilidad dentro de ese entorno \parencite{posix_shell, bash_best_practices}. Con ello, el alcance queda concentrado en una única implementación en Bash, en lugar de mantener además una versión paralela para otros sistemas operativos.

La robustez de los scripts se apoya en prácticas consolidadas del lenguaje: activar el modo estricto de ejecución (\texttt{set -euo pipefail}), entrecomillar adecuadamente las expansiones, emplear arrays y funciones en lugar de cadenas de comandos y verificar explícitamente los códigos de salida de cada instrucción \parencite{bash_best_practices}. Sobre esa base, el algoritmo incorpora el principio de idempotencia: una operación idempotente produce el mismo resultado sin efectos secundarios aunque se aplique repetidamente, de modo que re-ejecutar la automatización sobre una \vm{} ya configurada no reformatea discos ni duplica entradas en \texttt{/etc/fstab} \parencite{idempotency_concept, ansible_up_running}.

La calidad de ese código se verifica con dos herramientas complementarias. ShellCheck es un analizador estático que examina los scripts antes de ejecutarlos y detecta errores comunes y malas prácticas, como expansiones sin entrecomillar, comandos inexistentes o manejo deficiente de la salida, sin necesidad de correr el programa \parencite{shellcheck}. BATS (\textit{Bash Automated Testing System}), en cambio, ejecuta pruebas automatizadas sobre el comportamiento del script, incluida la re-ejecución idempotente sobre una \vm{} ya preparada \parencite{bats_testing}. De este modo, los defectos se detectan en etapa de desarrollo y no durante la configuración del sistema invitado.

\subsection{Testing de Scripts}

La verificación de los scripts combina análisis estático y pruebas funcionales. El análisis estático se realiza con ShellCheck sobre el código fuente, lo que permite identificar defectos potenciales en etapa de desarrollo y reduce el costo de corrección respecto a los fallos detectados durante la ejecución sobre una \vm{} \parencite{shellcheck}. Dicha verificación se ejecuta de forma previa a cualquier prueba funcional, de manera que solo los scripts que superan el linter avanzan a la siguiente etapa \parencite{bash_best_practices}.

Las pruebas funcionales se implementan con BATS, framework diseñado específicamente para scripts Bash, que organiza cada caso mediante bloques de configuración y limpieza (\textit{setup}/\textit{teardown}) y aserciones, y que produce salida en formato TAP, compatible con herramientas de integración continua \parencite{bats_testing}. En este trabajo, BATS se emplea para verificar los comportamientos críticos de la automatización: la preparación correcta de un disco nuevo, la detección de una configuración existente y, en particular, la re-ejecución sobre una \vm{} ya preparada, la cual debe finalizar con éxito sin modificar los datos persistidos \parencite{idempotency_concept}.
